\documentclass[12pt]{article}
\usepackage[utf8]{inputenc}
\usepackage[T1]{fontenc}
\usepackage{hyperref}
\title{Tomyang Language Resources for K-12 Curriculum Development}
\author{Compiled by Claude Sonnet 5 (claude-sonnet-5) for Andr\'as Kornai}
\date{}
\begin{document}
\maketitle

\section{Introduction}
Tomyang (Chongka) was ISO 639-3 tmx (Glottocode tomy1239), a code retired in 2007 and merged into Yamphu [ybi]. Glottolog lists it under ``Bookkeeping''. The only data collected was a Swadesh 100-word list, which showed 82 percent lexical similarity with Yamphu. Wikipedia's Yamphu article describes Tomyang as a dialect spoken in Nepal by about 20 people.

\section{Primary Educational Resources}
None found.

\section{Linguistic and Reference Materials}
\begin{itemize}
\item Glottolog entry for Tomyang: \url{https://glottolog.org/resource/languoid/id/tomy1239}
\item Wikipedia, Yamphu language, which mentions Tomyang: \url{https://en.wikipedia.org/wiki/Yamphu_language}
\end{itemize}

\section{Religious and Community Resources}
None found.

\section{Digital and Audio Resources}
None found.

\section{Curriculum Development Context}
Tomyang is not a separate language in ISO 639-3, and it has no known educational resources. Any curriculum work would fall under Yamphu.
\end{document}